% !TeX program = xelatex
%
% IMPE LaTeX System
% Canonical Showcase
%
% Intended location:
%   _showcase/main.tex
%
% Suggested build:
%
%   xelatex -shell-escape main.tex
%   biber main
%   xelatex -shell-escape main.tex
%   xelatex -shell-escape main.tex
%
% Font declarations are intentionally completed in the preamble.  Most script
% families are loaded as explicit local interfaces only; automatic Unicode-
% range routing is enabled for a small representative set.  This avoids the
% current v0.1.3 interaction between overlapping automatic owners and explicit
% local overrides, while keeping the routing regression visible.
%
% The first style matrix is deferred to the following landscape page so that
% the surrounding portrait text may continue normally.  This relies on the
% tables feature loading afterpage.
%
% Mathematical font selection is left at the IMPE default.  The showcase does
% not force Libertinus Math, and typewriter text is reset to Latin Modern Mono.
%
% Egyptian Hieroglyphs remain intentionally omitted from the showcase.
% Khitan Small Script is included only as the current linear fallback;
% canonical stacked Khitan rendering is not yet claimed here.

\begin{filecontents*}[overwrite]{references.bib}
@book{knuth1984texbook,
  author    = {Donald E. Knuth},
  title     = {The TeXbook},
  date      = {1984},
  publisher = {Addison-Wesley}
}
\end{filecontents*}

\documentclass[11pt,twoside]{nextart}

\UseCitationStyle{APA}

\HeaderTitle{IMPE Canonical Showcase}
\HeaderStyle{running}

% IMPORTANT:
% Do not insert blank lines inside the registration argument blocks.

\UseTemplateSet{
  fonts = {
    libertinus,
    shanggu
  },
  features = {
    math,
    tables,
    image,
    lists_envs,
    citations,
    index,
    hyperlinks,
    headers
  }
}

\setmonofont{Latin Modern Mono}

% Explicit local interfaces.  These do not claim automatic Unicode ranges.
\UseFont{noto}[local]
\UseFont{times}[local]
\UseFont{gentium}[local]
\UseFont{charis}[local]
\UseFont{anatolian}[local]
\UseFont{bopomofo}[local]
\UseFont{cuneiform}[local]
\UseFont{italic}[local]
\UseFont{hungarian}[local]
\UseFont{runic}[local]
\UseFont{hindi}[local]
\UseFont{sanskrit}[local]
\UseFont{devanagari}[local]
\UseFont{brahmi}[local]
\UseFont{tibetan}[local]
\UseFont{arabic}[local]
\UseFont{urdu}[local]
\UseFont{aramaic}[local]
\UseFont{nabataean}[local]
\UseFont{hebrew}[local]
\UseFont{syriac}[local]
\UseFont{syriac_eastern}[local]
\UseFont{kharosthi}[local]
\UseFont{pahlavi_parthian}[local]
\UseFont{pahlavi_inscriptional}[local]
\UseFont{pahlavi_psalter}[local]
\UseFont{avestan}[local]
\UseFont{manichaean}[local]
\UseFont{phoenician}[local]
\UseFont{samaritan}[local]
\UseFont{sogdian}[local]
\UseFont{sogdian_old}[local]
\UseFont{chinese_simplified}[local]
\UseFont{chinese_traditional}[local]
\UseFont{japanese}[local]
\UseFont{wenjin}[local]
\UseFont{korean}[local]
\UseFont{tangut}[local]
\UseFont{khitan_small}[local]
\UseFont{mongolian}[local]
\UseFont{mongolian_baiti}[local]
\UseFont{manchu}[local]
\UseFont{segoe}[local]
\UseFont{thai}[local]
\UseFont{turkic}[local]
\UseFont{uyghur}[local]
\UseFont{vietnamese_quocngu}[local]
\UseFont{vietnamese_hannom}[local]

% Minimal automatic Unicode-range regression set.
\UseFont{armenian}
\UseFont{georgian}
\UseFont{coptic}
\UseFont{glagolitic}
\UseFont{tamil}

\addbibresource{references.bib}

\hypersetup{
  pdftitle  = {IMPE LaTeX System — Canonical Showcase},
  pdfauthor = {IMPE LaTeX System v0.1.3},
  pdfsubject = {Canonical multilingual and feature showcase}
}

\setcounter{tocdepth}{2}
\raggedbottom

\newcommand{\ShowcaseStyles}[2]{%
  #1{#2} &
  #1{\textbf{#2}} &
  #1{\textit{#2}} &
  #1{\textbf{\textit{#2}}} &
  #1{\textsf{#2}} &
  #1{\textsf{\textbf{#2}}} &
  #1{\textsf{\textit{#2}}} &
  #1{\textsf{\textbf{\textit{#2}}}}%
}

\newcommand{\StyleRow}[4]{%
  #1 &
  \texttt{\textbackslash#2} &
  \ShowcaseStyles{#3}{#4} \\
}

\newcommand{\StyleCategory}[1]{%
  \midrule
  \multicolumn{10}{l}{\textbf{#1}}\\
  \midrule
}

\title{IMPE \LaTeX~System}
\subtitle{Canonical Multilingual and Feature Showcase}
\author{Kelvin Yang}
\date{v0.1.3}

\begin{document}

\maketitle

\begin{abstract}
This document is the canonical visual specimen for the current IMPE \LaTeX~System public interface. It combines a local-family style audit, a deliberately
small Unicode-range routing regression, dense multilingual paragraph
composition with explicit local overrides, right-to-left text, vertical
writing, generic complex shaping, CJK and Han family selection, and the
principal document features provided by IMPE. It is intended both as a public
showcase and as a release-time visual regression specimen.\footnote{The
showcase is intentionally broader than the small regression examples under
\texttt{examples/}; those examples remain useful for isolating individual
failures.}
\end{abstract}

\tableofcontents



% ==========================================================================
\section{Local Font Families and Styles}
% ==========================================================================

This section presents IMPE's explicit local font commands and compares their
available serif, bold, italic, and sans-serif styles. The complete font-family
matrix appears on the following landscape page. Where a requested face is
unavailable, IMPE falls back according to the family configuration rather than
requiring special-case markup.

\afterpage{%
\begin{landscape}

\begin{TableLong}{
  @{}
  P{2.55cm}
  M{1.25cm}
  *{8}{M{1.80cm}}
  @{}
}
\toprule
Family &
Command &
Regular &
Bold &
Italic &
Bold italic &
Sans &
Sans bold &
Sans italic &
Sans bold italic
\\
\midrule
\endfirsthead

\toprule
Family &
Command &
Regular &
Bold &
Italic &
Bold italic &
Sans &
Sans bold &
Sans italic &
Sans bold italic
\\
\midrule
\endhead

\StyleCategory{Latin and general-purpose families}
\StyleRow{Noto}{NOT}{\NOT}{Abg fi fl}
\StyleRow{Times}{TIM}{\TIM}{Abg fi fl}
\StyleRow{Gentium Plus}{GEN}{\GEN}{Ārya ṛ ṃ ḥ}
\StyleRow{Charis SIL}{CHA}{\CHA}{Ārya ṛ ṃ ḥ}
\StyleRow{Libertinus}{LIB}{\LIB}{Abg fi fl}

\StyleCategory{European and ancient alphabetic scripts}
\StyleRow{Carian}{CA}{\CA}{𐊠𐊺𐊪}
\StyleRow{Coptic}{CO}{\CO}{ⲛⲟⲩⲧⲉ}
\StyleRow{Bopomofo}{ZY}{\ZY}{ㄓㄨˋ ㄧㄣ}
\StyleRow{Cuneiform}{CU}{\CU}{𒀭𒂗𒆠}
\StyleRow{Glagolitic}{GL}{\GL}{ⰀⰁⰂ}
\StyleRow{Old Italic}{OI}{\OI}{𐌀𐌁𐌂}
\StyleRow{Old Hungarian}{OH}{\OH}{𐲀𐲁𐲂}
\StyleRow{Runic}{RU}{\RU}{ᚠᚢᚦ}

\StyleCategory{South and Central Asian scripts}
\StyleRow{Armenian}{HY}{\HY}{Հայերեն}
\StyleRow{Hindi}{HI}{\HI}{हिन्दी}
\StyleRow{Sanskrit}{SA}{\SA}{संस्कृतम्}
\StyleRow{Devanagari}{DEV}{\DEV}{देवनागरी}
\StyleRow{Tamil}{TA}{\TA}{தமிழ்}
\StyleRow{Brahmi}{BR}{\BR}{𑀩𑁆𑀭𑀸𑀳𑁆𑀫𑀻}
\StyleRow{Georgian}{KA}{\KA}{ქართული}
\StyleRow{Tibetan}{TI}{\TI}{བསྒྲུབས་}

\StyleCategory{RTL and West Asian scripts}
\StyleRow{Arabic}{AR}{\AR}{العربية}
\StyleRow{Urdu}{UR}{\UR}{اُردُو}
\StyleRow{Imperial Aramaic}{IA}{\IA}{𐡀𐡓𐡌}
\StyleRow{Nabataean}{NB}{\NB}{𐢌𐢁𐢈}
\StyleRow{Hebrew}{HE}{\HE}{שלום}
\StyleRow{Syriac}{SY}{\SY}{ܫܠܡܐ}
\StyleRow{Eastern Syriac}{SYE}{\SYE}{ܫܠܡܐ}
\StyleRow{Kharoshthi}{KH}{\KH}{𐨐𐨿𐨪𐨆}

\StyleCategory{Iranian and other historical scripts}
\StyleRow{Inscriptional Parthian}{PAR}{\PAR}{𐭐𐭓𐭕}
\StyleRow{Inscriptional Pahlavi}{PAH}{\PAH}{𐭯𐭠𐭧}
\StyleRow{Psalter Pahlavi}{PSP}{\PSP}{𐮈𐮎}
\StyleRow{Avestan}{AV}{\AV}{𐬀𐬬𐬆𐬯}
\StyleRow{Manichaean}{MA}{\MA}{𐫀𐫁𐫂}
\StyleRow{Phoenician}{PH}{\PH}{𐤀𐤁𐤂}
\StyleRow{Samaritan}{SM}{\SM}{ࠔࠌࠅࠌ}
\StyleRow{Sogdian}{SG}{\SG}{𐼰𐼱𐼲}
\StyleRow{Old Sogdian}{SGO}{\SGO}{𐼀𐼁𐼂}

\StyleCategory{CJK and East Asian families}
\StyleRow{Noto SC}{SC}{\SC}{汉字骨直}
\StyleRow{Noto TC}{TC}{\TC}{漢字骨直}
\StyleRow{Japanese}{JP}{\JP}{日本語漢字}
\StyleRow{WenJin}{WJ}{\WJ}{中國𠀀𰀀}
\StyleRow{Korean}{KR}{\KR}{한국어漢字}
\StyleRow{Tangut}{TG}{\TG}{𗼇𘝞𗫂}

\StyleCategory{Mongolian-family and other Asian scripts}
\StyleRow{Khitan Small Script (linear)}{KHS}{\KHS}{𘬜𘭢𘮒}
\StyleRow{Mongolian}{MO}{\MO}{ᠮᠣᠩᠭᠣᠯ}
\StyleRow{Mongolian Baiti}{MOb}{\MOb}{ᠮᠣᠩᠭᠣᠯ}
\StyleRow{Manchu}{MC}{\MC}{ᠮᠠᠨᠵᡠ}
\StyleRow{Segoe Historic}{SEG}{\SEG}{𑀩𑁆𑀭𑀸𑀳𑁆𑀫𑀻}
\StyleRow{Thai}{TH}{\TH}{ภาษาไทย}
\StyleRow{Old Turkic}{OT}{\OT}{𐱅𐰇𐰼𐰚}
\StyleRow{Old Uyghur}{UY}{\UY}{𐽳𐽶𐽲𐽳𐽾}

\StyleCategory{Vietnamese}
\StyleRow{Quốc Ngữ}{VI}{\VI}{Tiếng Việt}
\StyleRow{Hán-Nôm}{HN}{\HN}{𡨸喃𠊛𡗶}

\bottomrule
\end{TableLong}

\end{landscape}
}


% ==========================================================================
\section{Multilingual Routing and Mixed-Script Composition}
% ==========================================================================

\Term[sort=multilingual composition]{Multilingual composition}
[combined routing and local overrides] is tested here in two deliberately
different ways. The main specimen is a dense paragraph in which explicit local
families, right-to-left runs, complex shaping, shared Han choices, and inline
vertical writing are allowed to interact. A later compact regression enables
automatic Unicode-range ownership for only a small set of ordinary scripts.

\subsection{Mixed-Script Paragraph}

The following specimen is intentionally one \TeX~paragraph. It contains no
English transition sentences: every script that appears contributes a complete
sentence or clause. Armenian, Georgian, Coptic, Glagolitic, and Tamil are
left as raw Unicode to exercise the small automatic-routing set; all other
script families in the paragraph use explicit local overrides.

Հայերեն այս ամբողջական նախադասությունը գրված է նույն բազմալեզու պարբերության մեջ։
ეს სრული ქართული წინადადება იმავე მრავალენოვან აბზაცში სხვა დამწერლობებთან ერთად გვხვდება.
Ⲟⲩⲟϩ ⲡⲓⲗⲟⲅⲟⲥ ⲁϥϣⲱⲡⲓ ⲛ̀ⲟⲩⲥⲁⲣⲝ ⲟⲩⲟϩ ⲁϥϣⲱⲡⲓ ⲛ̀ϧⲏⲧⲉⲛ.
Ⰲⱏ ⱀⰰⱍⰰⰾⱑ ⰱⱑⰰⱎⰵ ⱄⰾⱁⰲⱁ, ⰺ ⱄⰾⱁⰲⱁ ⰱⱑⰰⱎⰵ ⰽⱏ ⰱⱁⰳⱆ.
\SA{सर्वे भवन्तु सुखिनः, सर्वे सन्तु निरामयाः।}
தமிழில் எழுதப்பட்ட இந்த முழு வாக்கியம் அதே பலமொழிப் பத்திக்குள் தொடர்ந்து வருகிறது.
\TI{ཤེས་རབ་དང་སྙིང་རྗེ་ནི་མིའི་ཆོས་ཀྱི་རྩ་བ་ཡིན།}
\AR{هذه جملة عربية كاملة تظهر في وسط الفقرة نفسها وتختبر الانتقال إلى الكتابة من اليمين إلى اليسار.}
\UR{یہ ایک مکمل اردو جملہ ہے جو اسی کثیر لسانی پیراگراف کے اندر نستعلیق رسم الخط میں لکھی گئی ہے۔}
\HE{זהו משפט עברי שלם שממשיך באותה פסקה לאחר הטקסט הערבי והאורדו.}
\SY{ܗܢܐ ܦܬܓܡܐ ܣܘܪܝܝܐ ܟܠܗ ܟܬܝܒ ܒܓܘ ܗܕܐ ܦܪܓܪܦܐ ܥܡ ܠܫܢ̈ܐ ܐܚܪ̈ܢܐ.}
\SYE{ܗܢܐ ܦܬܓܡܐ ܐܚܪܢܐ ܟܬܝܒ ܒܟܬܒܐ ܡܕܢܚܝܐ ܘܡܫܬܠܚ ܒܓܘ ܗܕܐ ܦܪܓܪܦܐ.}
\SEG{𑀲𑀢𑁆𑀬𑀁 𑀯𑀤, 𑀥𑀭𑁆𑀫𑀁 𑀘𑀭.}
\JP{日本語のこの文章では、ひらがなとカタカナと漢字が同じ文の中に自然に現れます。}
\KR{한국어의 이 문장에서는 한글과 漢字가 같은 문장 안에 함께 나타납니다.}
\TC{這個中文完整句子使用明確指定的繁體中文字體，並與前後不同文種保持在同一個段落中。}
\MOv{ᠮᠣᠩᠭᠣᠯ ᠬᠡᠯᠡ ᠪᠣᠯ ᠮᠣᠩᠭᠣᠯᠴᠤᠳ ‍ᠤᠨ ᠬᠡᠯᠡ ᠮᠥᠨ᠂ ᠮᠣᠩᠭᠣᠯ ᠪᠢᠴᠢᠭ ᠨᠢ ᠪᠣᠰᠤᠭ᠎ᠠ ᠴᠢᠭᠯᠡᠯ ᠢᠶᠡᠷ ᠪᠢᠴᠢᠭᠳᠡᠨ᠎ᠡ᠃}
\JP{縦書きのモンゴル文字の直後にも日本語の完全な文章が続き、通常の横書き状態が回復したことを確認します。}
\MCv{ᠪᡳ ᠮᠠᠨᠵᡠ ᡤᡳᠰᡠᠨ ᠪᡝ ᡨᠠᠴᡳᠮᠪᡳ᠂ ᡝᡵᡝ ᠪᡳᡨᡥᡝ ᡩᡝ ᠪᠣᠰᠣ ᠠᡵᠠᠮᠪᡳ᠉}
\KR{두 번째 세로쓰기 구간 뒤에도 한국어의 완전한 문장이 이어지며 문단의 방향과 글꼴 상태가 정상적으로 복원되는지를 확인합니다.}
\TC{最後這個中文完整句子仍然留在同一個段落裡，用來確認兩次豎排 local override、數次 RTL 切換與複雜 shaping 之後，普通橫排行文與局部字體狀態全部恢復正常。}

\subsection{Minimal Automatic Unicode-Range Regression}

\Term[sort=Unicode range routing]{Unicode-range routing}
[automatic font ownership] is exercised here only on a small representative
set. This keeps the v0.1.3 showcase focused on routing correctness without
turning the entire document into a large interchar transition graph. The
specimen itself contains raw Unicode only.


Հայերեն այս ամբողջական նախադասությունը գրված է առանց տեղական տառատեսակի հրամանի։
ეს სრული ქართული წინადადება ავტომატური Unicode დიაპაზონის მარშრუტიზაციას ამოწმებს.
Ⲟⲩⲟϩ ⲡⲓⲗⲟⲅⲟⲥ ⲁϥϣⲱⲡⲓ ⲛ̀ⲟⲩⲥⲁⲣⲝ ⲟⲩⲟϩ ⲁϥϣⲱⲡⲓ ⲛ̀ϧⲏⲧⲉⲛ.
Ⰲⱏ ⱀⰰⱍⰰⰾⱑ ⰱⱑⰰⱎⰵ ⱄⰾⱁⰲⱁ, ⰺ ⱄⰾⱁⰲⱁ ⰱⱑⰰⱎⰵ ⰽⱏ ⰱⱁⰳⱆ.
தமிழில் எழுதப்பட்ட இந்த முழு வாக்கியம் தனி local font command இன்றி Unicode range owner மூலம் அமைக்கப்படுகிறது.


% ==========================================================================
\section{Right-to-Left Scripts}
% ==========================================================================

\Term[sort=right-to-left text]{Right-to-left text}
[directional behavior] requires more than choosing a glyph source. This
section therefore uses explicit local families throughout. Each specimen is
long enough to form a genuine paragraph, allowing line breaking, inter-word
spacing, punctuation, digits, bidirectional interaction, and OpenType shaping
to be inspected without enabling overlapping automatic Arabic, Hebrew, or
Syriac range owners.

\subsection{Arabic}

\AR{%
اَلْعَرَبِيَّةُ لُغَةٌ ذاتُ تاريخٍ طويلٍ، وَهِيَ تُكْتَبُ مِنَ الْيَمِينِ إِلَى الْيَسَارِ، وَتَتَّصِلُ حُرُوفُهَا بِأَشْكَالٍ مُخْتَلِفَةٍ بِحَسَبِ مَوْضِعِهَا فِي الْكَلِمَةِ. يَحْتَوِي هٰذَا النَّصُّ عَلَى جُمَلٍ طَوِيلَةٍ وَعَلَامَاتِ تَرْقِيمٍ مُتَنَوِّعَةٍ، وَتَتَغَيَّرُ أَطْوَالُ الْكَلِمَاتِ وَالْجُمَلِ حَتَّى يَظْهَرَ سُلُوكُ الْفَقْرَةِ عِنْدَ نِهَايَةِ السَّطْرِ. وَعِنْدَمَا يَطُولُ السَّطْرُ يَجِبُ أَنْ يَنْتَقِلَ النَّصُّ إِلَى السَّطْرِ التَّالِي بِصُورَةٍ طَبِيعِيَّةٍ مِنْ غَيْرِ أَنْ تَنْقَطِعَ أَشْكَالُ الْحُرُوفِ أَوْ تَضْطَرِبَ الْمَسَافَاتُ بَيْنَ الْكَلِمَاتِ. وَيُسْتَعْمَلُ هٰذَا الْمَقْطَعُ لِلتَّأَكُّدِ مِنْ أَنَّ الْفَقْرَةَ الْعَرَبِيَّةَ الْكَامِلَةَ تَبْقَى مُسْتَقِرَّةً دَاخِلَ وَثِيقَةٍ مُتَعَدِّدَةِ اللُّغَاتِ.
\par
}

Inline Arabic inside an English sentence: \AR{العربية والكتابة في سطر واحد مع علامات الترقيم} before returning to ordinary Latin text.

\subsection{Urdu}

\UR{%
اُردُو ایک ایسی زبان ہے جس کی تحریری روایت میں نستعلیق رسم الخط کو خاص اہمیت حاصل ہے۔ اس رسم الخط میں حروف صرف ایک دوسرے سے ملتے ہی نہیں بلکہ الفاظ کی مجموعی ساخت بھی سطر کی بصری صورت پر اثر ڈالتی ہے، اس لیے طویل عبارت مختصر نمونے سے کہیں زیادہ مفید آزمائش فراہم کرتی ہے۔ اس پیراگراف میں مختلف لمبائی کے الفاظ، مختلف رموزِ اوقاف، اور کئی مسلسل جملے شامل ہیں تاکہ دائیں سے بائیں بہاؤ، لفظی فاصلوں، اور سطر کے اختتام پر عبارت کی تقسیم کو واضح طور پر دیکھا جا سکے۔ جب عبارت ایک سطر میں پوری نہ آئے تو اگلی سطر میں منتقلی قدرتی ہونی چاہیے، الفاظ کی ترتیب برقرار رہنی چاہیے، اور نستعلیق کی تشکیل کسی غلط فونٹ یا غیر متعلقہ عربی طرز میں تبدیل نہیں ہونی چاہیے۔ یہ آخری جملہ خاص طور پر اس بات کی جانچ کرتا ہے کہ ایک مکمل اُردُو پیراگراف کئی سطروں میں تقسیم ہونے کے بعد بھی اپنی مقامی نستعلیق فیملی اور درست سمت برقرار رکھتا ہے۔
\par
}

\subsection{Hebrew}

\HE{%
העברית נכתבת מימין לשמאל, ולכן פסקה ארוכה מאפשרת לבדוק לא רק את צורת האותיות אלא גם את התנהגות השורות, הרווחים וסימני הפיסוק. בטקסט הזה מופיעים משפטים באורכים שונים כדי לבדוק כיצד המערכת מתמודדת עם רצף עברי ארוך בתוך פסקה רגילה. כאשר השורה נעשית ארוכה מדי, המעבר לשורה הבאה צריך להתרחש במקום טבעי בלי לשנות את סדר המילים ובלי לפגוע במיקום הנקודות, הפסיקים והסוגריים. אפשר גם לכתוב מילים מנוקדות כמו שָׁלוֹם, עִבְרִית וְכְתִיבָה בתוך אותה פסקה, וכך לבדוק שהניקוד נשאר מחובר לאות הנכונה גם לאחר שבירת שורה. המשפט האחרון נועד לוודא שהפסקה כולה נשארת יציבה בתוך מסמך רב־לשוני ולאחריה חוזר הטקסט הסובב להתנהגות הרגילה.
\par
}

\subsection{Syriac}

Western / general Syriac:

\SY{%
ܫܠܡܐ ܥܠܝܟܘܢ. ܗܢܐ ܟܬܒܐ ܟܬܝܒ ܒܠܫܢܐ ܣܘܪܝܝܐ، ܘܐܝܬ ܒܗ ܡ̈ܠܐ ܣܓܝ̈ܐܬܐ ܘܦܬܓ̈ܡܐ ܕܐܪܝܟܝܢ ܡܢ ܚܕ ܫܛܪܐ. ܟܕ ܡܬܡܠܐ ܫܛܪܐ، ܒܥܐ ܟܬܒܐ ܕܢܚܘܬ ܠܫܛܪܐ ܐܚܪܢܐ ܒܠܐ ܕܢܬܚܠܦ ܛܟܣܐ ܕܡ̈ܠܐ ܘܒܠܐ ܕܢܬܦܣܩ ܚܘܒܪܐ ܕܐܬܘ̈ܬܐ. ܗܟܢܐ ܡܨܝܢܢ ܕܢܒܕܘܩ ܐܝܟܢܐ ܡܬܕܒܪ ܟܬܒܐ ܟܕ ܐܘܪܟܗ ܪܒ ܘܟܕ ܐܝܬ ܒܗ ܪܘܚ̈ܩܐ ܘܦܘ̈ܣܩܐ ܣܓܝ̈ܐܐ. ܗܢܐ ܦܬܓܡܐ ܐܚܪܝܐ ܡܘܣܦ ܐܘܪܟܐ ܠܦܪܓܪܦܐ ܡܛܠ ܕܢܚܙܐ ܣܓܝ̈ܐܐ ܫܛܪ̈ܐ ܘܢܒܕܘܩ ܪܘܚܩܐ ܘܛܟܣܐ ܕܟܬܒܐ.
\par
}

Eastern Syriac stylistic family:

\SYE{%
ܫܠܡܐ ܥܠܝܟܘܢ. ܗܢܐ ܟܬܒܐ ܟܬܝܒ ܒܠܫܢܐ ܣܘܪܝܝܐ، ܘܐܝܬ ܒܗ ܡ̈ܠܐ ܣܓܝ̈ܐܬܐ ܘܦܬܓ̈ܡܐ ܕܐܪܝܟܝܢ ܡܢ ܚܕ ܫܛܪܐ. ܟܕ ܡܬܡܠܐ ܫܛܪܐ، ܒܥܐ ܟܬܒܐ ܕܢܚܘܬ ܠܫܛܪܐ ܐܚܪܢܐ ܒܠܐ ܕܢܬܚܠܦ ܛܟܣܐ ܕܡ̈ܠܐ ܘܒܠܐ ܕܢܬܦܣܩ ܚܘܒܪܐ ܕܐܬܘ̈ܬܐ. ܗܟܢܐ ܡܨܝܢܢ ܕܢܒܕܘܩ ܐܝܟܢܐ ܡܬܕܒܪ ܟܬܒܐ ܟܕ ܐܘܪܟܗ ܪܒ ܘܟܕ ܐܝܬ ܒܗ ܪܘܚ̈ܩܐ ܘܦܘ̈ܣܩܐ ܣܓܝ̈ܐܐ. ܗܢܐ ܦܬܓܡܐ ܐܚܪܝܐ ܡܘܣܦ ܐܘܪܟܐ ܠܦܪܓܪܦܐ ܡܛܠ ܕܢܚܙܐ ܣܓܝ̈ܐܐ ܫܛܪ̈ܐ ܘܢܒܕܘܩ ܪܘܚܩܐ ܘܛܟܣܐ ܕܟܬܒܐ.
\par
}

% ==========================================================================
\section{Vertical Writing}
% ==========================================================================

\Term[sort=vertical writing]{Vertical writing}
[horizontal-layout-compatible vertical embedding] is currently implemented
as a compatibility mode inside an otherwise horizontal document layout.
IMPE does not yet provide a fully vertical page or paragraph flow. Instead,
the text is divided at word boundaries, and each word is typeset vertically
while the resulting word-sized units continue to participate in ordinary
horizontal paragraph layout. Spaces therefore become natural break
opportunities between vertical words.

This design keeps genuinely vertical glyph orientation while allowing the
text to wrap within normal horizontal pages, with relatively compact line
spacing and without requiring a separate vertical document model. The
specimens below are deliberately long enough to test this compatibility
strategy under realistic paragraph conditions: word segmentation, inter-word
spacing, line breaking, line spacing, punctuation, ordering, and recovery of
the surrounding horizontal layout.

\subsection{Mongolian}

\noindent
\begin{minipage}[t]{0.47\linewidth}
\raggedright
\textbf{Horizontal}

\medskip

\MO{ᠮᠣᠩᠭᠣᠯ ᠬᠡᠯᠡ ᠪᠣᠯ ᠮᠣᠩᠭᠣᠯᠴᠤᠳ ‍ᠤᠨ ᠬᠡᠯᠡ ᠮᠥᠨ᠃ ᠮᠣᠩᠭᠣᠯ ᠪᠢᠴᠢᠭ ᠨᠢ ᠪᠣᠰᠤᠭ᠎ᠠ ᠴᠢᠭᠯᠡᠯ ᠢᠶᠡᠷ ᠪᠢᠴᠢᠭᠳᠡᠨ᠎ᠡ᠂ ᠦᠭᠡ ᠪᠦᠷᠢ ᠨᠢ ᠬᠣᠭᠣᠷᠣᠨᠳᠣᠭ᠎ᠠ ᠵᠠᠢ ᠲᠠᠢ ᠪᠠᠢᠳᠠᠭ᠃ ᠤᠷᠲᠤ ᠦᠭᠡ ᠪᠣᠯᠣᠨ ᠣᠯᠠᠨ ᠦᠭᠡ ᠨᠢ ᠨᠢᠭᠡ ᠮᠥᠷ ᠲᠦ ᠪᠠᠭᠲᠠᠬᠤ ᠦᠭᠡᠢ ᠦᠶ᠎ᠡ ᠳᠦ ᠳᠠᠷᠠᠭᠠᠴᠢ ᠶᠢᠨ ᠮᠥᠷ ᠲᠦ ᠰᠢᠯᠵᠢᠨ᠎ᠡ᠃}
\end{minipage}
\hfill
\begin{minipage}[t]{0.47\linewidth}
\raggedright
\textbf{Vertical}

\medskip

\MOv{ᠮᠣᠩᠭᠣᠯ ᠬᠡᠯᠡ ᠪᠣᠯ ᠮᠣᠩᠭᠣᠯᠴᠤᠳ ‍ᠤᠨ ᠬᠡᠯᠡ ᠮᠥᠨ᠃ ᠮᠣᠩᠭᠣᠯ ᠪᠢᠴᠢᠭ ᠨᠢ ᠪᠣᠰᠤᠭ᠎ᠠ ᠴᠢᠭᠯᠡᠯ ᠢᠶᠡᠷ ᠪᠢᠴᠢᠭᠳᠡᠨ᠎ᠡ᠂ ᠦᠭᠡ ᠪᠦᠷᠢ ᠨᠢ ᠬᠣᠭᠣᠷᠣᠨᠳᠣᠭ᠎ᠠ ᠵᠠᠢ ᠲᠠᠢ ᠪᠠᠢᠳᠠᠭ᠃ ᠤᠷᠲᠤ ᠦᠭᠡ ᠪᠣᠯᠣᠨ ᠣᠯᠠᠨ ᠦᠭᠡ ᠨᠢ ᠨᠢᠭᠡ ᠮᠥᠷ ᠲᠦ ᠪᠠᠭᠲᠠᠬᠤ ᠦᠭᠡᠢ ᠦᠶ᠎ᠡ ᠳᠦ ᠳᠠᠷᠠᠭᠠᠴᠢ ᠶᠢᠨ ᠮᠥᠷ ᠲᠦ ᠰᠢᠯᠵᠢᠨ᠎ᠡ᠃}
\end{minipage}

\bigskip

\subsection{Manchu}

\noindent
\begin{minipage}[t]{0.47\linewidth}
\raggedright
\textbf{Horizontal}

\medskip

\MC{ᠮᠠᠨᠵᡠ ᡤᡳᠰᡠᠨ ᠪᡝ ᠮᠠᠨᠵᡠ ᠪᡳᡨᡥᡝ ‍ᡳ ᠠᡵᠠᠮᠪᡳ᠃ ᠪᡳ ᡝᡵᡝ ᠪᡳᡨᡥᡝ ᡩᡝ ᠯᠣᠨᡤᡤᠣ ᡤᡳᠰᡠᠨ ᠪᡝ ᠠᡵᠠᠮᠪᡳ᠂ ᡨᡝᡵᡝᠴᡳ ᡤᡳᠰᡠᠨ ᡨᠠᡨᠠᠮᡝ ᡤᡝᠨᡝᠮᠪᡳ᠃ ᡝᡵᡝ ᡳᠴᡳ ᠠᠮᠪᠠ ᠪᡳᡨᡥᡝ ᠪᡝ ᠠᡵᠠᡥᠠ ᠮᠠᠩᡤᡳ ᠮᡠᡵᡠ ᠪᡝ ᡥᡝᠴᡝᡥᡝ᠂ ᡤᡳᠰᡠᠨ ᠪᡝ ᠨᠠᡩᠠᠨ ᠪᠠ ᡩᡝ ᡶᠠᡵᠠᡥᠠ ᠪᡝ ᡳᠯᡝᡨᡠᠯᡝᠮᠪᡳ᠃}
\end{minipage}
\hfill
\begin{minipage}[t]{0.47\linewidth}
\raggedright
\textbf{Vertical}

\medskip

\MCv{ᠮᠠᠨᠵᡠ ᡤᡳᠰᡠᠨ ᠪᡝ ᠮᠠᠨᠵᡠ ᠪᡳᡨᡥᡝ ‍ᡳ ᠠᡵᠠᠮᠪᡳ᠃ ᠪᡳ ᡝᡵᡝ ᠪᡳᡨᡥᡝ ᡩᡝ ᠯᠣᠨᡤᡤᠣ ᡤᡳᠰᡠᠨ ᠪᡝ ᠠᡵᠠᠮᠪᡳ᠂ ᡨᡝᡵᡝᠴᡳ ᡤᡳᠰᡠᠨ ᡨᠠᡨᠠᠮᡝ ᡤᡝᠨᡝᠮᠪᡳ᠃ ᡝᡵᡝ ᡳᠴᡳ ᠠᠮᠪᠠ ᠪᡳᡨᡥᡝ ᠪᡝ ᠠᡵᠠᡥᠠ ᠮᠠᠩᡤᡳ ᠮᡠᡵᡠ ᠪᡝ ᡥᡝᠴᡝᡥᡝ᠂ ᡤᡳᠰᡠᠨ ᠪᡝ ᠨᠠᡩᠠᠨ ᠪᠠ ᡩᡝ ᡶᠠᡵᠠᡥᠠ ᠪᡝ ᡳᠯᡝᡨᡠᠯᡝᠮᠪᡳ᠃}
\end{minipage}

\bigskip

\subsection{Old Uyghur}

The Old Uyghur sample is deliberately extended so that several successive
vertical columns must be produced. It is used as a layout stress string rather
than as a continuous philological quotation.

\noindent
\begin{minipage}[t]{0.47\linewidth}
\raggedright
\textbf{Horizontal}

\medskip

\UY{𐽳𐽶𐽲𐽳𐽾 𐽰𐽶𐾁𐽶 𐽳𐽶𐽲𐽳𐽾 𐽰𐽶𐾁𐽶 𐽳𐽶𐽲𐽳𐽾 𐽰𐽶𐾁𐽶 𐽳𐽶𐽲𐽳𐽾 𐽰𐽶𐾁𐽶 𐽳𐽶𐽲𐽳𐽾 𐽰𐽶𐾁𐽶 𐽳𐽶𐽲𐽳𐽾 𐽰𐽶𐾁𐽶 𐽳𐽶𐽲𐽳𐽾 𐽰𐽶𐾁𐽶 𐽳𐽶𐽲𐽳𐽾 𐽰𐽶𐾁𐽶 𐽳𐽶𐽲𐽳𐽾 𐽰𐽶𐾁𐽶 𐽳𐽶𐽲𐽳𐽾 𐽰𐽶𐾁𐽶 𐽳𐽶𐽲𐽳𐽾 𐽰𐽶𐾁𐽶 𐽳𐽶𐽲𐽳𐽾 𐽰𐽶𐾁𐽶 𐽳𐽶𐽲𐽳𐽾 𐽰𐽶𐾁𐽶 𐽳𐽶𐽲𐽳𐽾 𐽰𐽶𐾁𐽶 𐽳𐽶𐽲𐽳𐽾 𐽰𐽶𐾁𐽶 𐽳𐽶𐽲𐽳𐽾 𐽰𐽶𐾁𐽶}
\end{minipage}
\hfill
\begin{minipage}[t]{0.47\linewidth}
\raggedright
\textbf{Vertical}

\medskip

\UYv{𐽳𐽶𐽲𐽳𐽾 𐽰𐽶𐾁𐽶 𐽳𐽶𐽲𐽳𐽾 𐽰𐽶𐾁𐽶 𐽳𐽶𐽲𐽳𐽾 𐽰𐽶𐾁𐽶 𐽳𐽶𐽲𐽳𐽾 𐽰𐽶𐾁𐽶 𐽳𐽶𐽲𐽳𐽾 𐽰𐽶𐾁𐽶 𐽳𐽶𐽲𐽳𐽾 𐽰𐽶𐾁𐽶 𐽳𐽶𐽲𐽳𐽾 𐽰𐽶𐾁𐽶 𐽳𐽶𐽲𐽳𐽾 𐽰𐽶𐾁𐽶 𐽳𐽶𐽲𐽳𐽾 𐽰𐽶𐾁𐽶 𐽳𐽶𐽲𐽳𐽾 𐽰𐽶𐾁𐽶 𐽳𐽶𐽲𐽳𐽾 𐽰𐽶𐾁𐽶 𐽳𐽶𐽲𐽳𐽾 𐽰𐽶𐾁𐽶 𐽳𐽶𐽲𐽳𐽾 𐽰𐽶𐾁𐽶 𐽳𐽶𐽲𐽳𐽾 𐽰𐽶𐾁𐽶 𐽳𐽶𐽲𐽳𐽾 𐽰𐽶𐾁𐽶 𐽳𐽶𐽲𐽳𐽾 𐽰𐽶𐾁𐽶}
\end{minipage}


% ==========================================================================
\section{Complex Shaping and Stacked Sequences}
% ==========================================================================

\Term[sort=OpenType shaping]{OpenType shaping}
[generic complex shaping] is tested here through explicit local families.
This deliberately separates shaping from Unicode-range ownership: if a local
specimen is wrong, the failure belongs to the selected family or shaping path
rather than to an automatic range transition. Egyptian Hieroglyphs remain
outside this showcase, and Khitan Small Script is recorded only in its current
linear fallback form.

\subsection{Devanagari}

The Sanskrit local family exercises conjuncts, virama behavior, pre-base and
post-base vowel signs, and longer clusters without enabling a document-wide
Devanagari owner.

\begin{TableBookX}
  {L L}
  {Devanagari shaping specimens}
  {tab:deva-shaping}
Sequence & Specimen \\
\midrule
kṣa & \SA{क्ष} \\
jña & \SA{ज्ञ} \\
śra & \SA{श्र} \\
stra & \SA{स्त्र} \\
ddhya & \SA{द्ध्य} \\
dharma & \SA{धर्म} \\
dharmakṣetre & \SA{धर्मक्षेत्रे} \\
saṃskṛtam & \SA{संस्कृतम्} \\
oṃ & \SA{ॐ} \\
\end{TableBookX}

A longer local specimen:

\begin{ExampleBlock}
\SA{धर्मक्षेत्रे कुरुक्षेत्रे समवेता युयुत्सवः।
संस्कृतवाक्यानि देवनागर्या लिखितानि सन्ति।}
\end{ExampleBlock}

\subsection{Tibetan}

Tibetan provides a visible example of generic vertical stacking inside one
local shaping run.

\begin{TableBookX}
  {L L}
  {Tibetan stack specimens}
  {tab:tibetan-shaping}
Type & Specimen \\
\midrule
stack & \TI{བསྒྲུབས་} \\
stack & \TI{བརྒྱུད་} \\
stack & \TI{བསྟན་} \\
stack & \TI{རྒྱལ་} \\
stack & \TI{དབྱངས་} \\
\end{TableBookX}

A longer local paragraph exercises Tibetan shaping and line breaking under
ordinary document width:

\TI{བོད་ཡིག་ནི་བོད་ཀྱི་སྐད་ཡིག་འབྲི་བའི་ཡི་གེའི་མ་ལག་ཅིག་ཡིན་ཞིང་། ཚིག་དང་ཚིག་གི་བར་ལ་ཚེག་གིས་མཚམས་འབྱེད་པ་དང་། ཡི་གེ་མང་པོ་ཞིག་མགོ་ཅན་དང་འདོགས་ཅན་དུ་སྦྱར་ནས་ཚེག་བར་གཅིག་གི་ནང་དུ་སྤུངས་ཏེ་འབྲི་དགོས།དཔེར་ན་བསྒྲུབས་པ་དང་། བརྒྱུད་པ། བསྟན་པ། རྒྱལ་ཁབ། དབྱངས་ཅན་ལྟ་བུའི་ཚིག་ནང་དུ་ཡི་གེ་མང་པོ་སྦྱར་ཡོད་པས། ཡིག་གཟུགས་ཀྱི་སྤུངས་སྟངས་དང་ཚེག་བར་གྱི་བར་ཐག་གཉིས་ཀ་ཡང་དག་པར་སྟོན་དགོས། ཚིག་ཕྲེང་རིང་པོ་ཞིག་ཤོག་ངོས་ཀྱི་ཞེང་ཚད་ལས་བརྒལ་བའི་སྐབས་སུ། ཚེག་གི་རྗེས་སུ་རང་བཞིན་གྱིས་ཕྲེང་བ་གསར་པར་འགྲོ་ཐུབ་པ་དང་། སྔ་ཕྱིའི་ཡི་གེའི་སྤུངས་སྟངས་དང་ཚིག་གི་གོ་རིམ་མི་འཁྲུག་པ་གལ་ཆེ། དེའི་ཕྱིར་ཚིག་ཕྲེང་རིང་པོ་འདིས་བོད་ཡིག་གི་གཟུགས་སྦྱོར་དང་། ཚེག་བར། ཕྲེང་བར། ཤད་དང་ཉིས་ཤད་ཀྱི་གནས་སྟངས། དེ་བཞིན་དུ་ཕྲེང་བ་བརྗེ་བའི་སྐབས་ཀྱི་སྤྱིའི་སྒྲིག་སྟངས་བཅས་ལ་ཚོད་ལྟ་བྱེད་ཐུབ།།}

\subsection{Brahmi, Kharoshthi, and Thai}

Brahmi is provided through two local families. The redistributable Noto
family provides broad Brahmi coverage, but in the current IMPE/\XeTeX~setup
it does not produce the expected stacked forms for all conjunct sequences.
Segoe UI Historic is therefore exposed through \texttt{\textbackslash SEG}
as an alternative reference family and produces the intended stacking in
these cases. Because Segoe UI Historic is a Microsoft Windows font, its font
binary is not distributed with IMPE; the local command is available only
when the user supplies or already has access to the font.

\begin{TableBookX}
  {L L}
  {Additional generic-shaping specimens}
  {tab:generic-shaping}
Script / family & Specimen \\
\midrule
Brahmi / Noto &
  \BR{𑀲𑀢𑁆𑀬𑀁 𑀯𑀤, 𑀥𑀭𑁆𑀫𑀁 𑀘𑀭.} \\
Brahmi / Segoe Historic &
  \SEG{𑀲𑀢𑁆𑀬𑀁 𑀯𑀤, 𑀥𑀭𑁆𑀫𑀁 𑀘𑀭.} \\
Kharoshthi &
  \KH{𐨢𐨿𐨪𐨨 𐨐𐨿𐨪𐨨 𐨯𐨿𐨪𐨬 𐨡𐨅𐨬 𐨣𐨨 𐨐𐨁 𐨐𐨂 𐨐𐨅 𐨐𐨆} \\
Thai &
  \TH{ภาษาไทย กำ กิ กี้ เก้า} \\
\end{TableBookX}

\subsection{Khitan Small Script: current linear fallback}

Khitan Small Script is already exposed through IMPE's public font interface,
while its canonical stacked layout is not yet used in this showcase. The
following specimen is taken from the Khitan Small Script epitaph of Xiao
Hudujin and is deliberately long enough to exercise the present linear
fallback across continuous text, cluster boundaries, spacing, and paragraph
line breaking.

\KHS{𘲇 𘰺𘱹𘰻 𘱿𘱤 𘱚𘮒𘲦 𘯣𘰴 𘲤𘭂𘮢𘲫𘲦 𘰺𘱄𘱯𘱸 𘬥𘲜𘰭 𘲜𘬚 𘭃𘱚𘱸 𘭞𘱄𘮽𘰕 𘲽𘭽𘭢 𘲽𘲑 𘰲 𘲅𘯺𘬐 𘲽𘱚𘲦 𘲚𘰥 𘮯𘮠 𘰷𘯀 𘱚𘱕 𘮧𘯶𘮱𘱄𘲫 𘬜𘭪𘲚𘱪 𘭩 𘱚𘱛 𘰭𘲚 𘮵𘯺𘲑 𘲽𘰥𘱦 𘮽𘯛 𘬜𘰆 𘰷𘱤 𘲽𘱦 𘭃𘭽𘲽 𘱤𘰣 𘱛𘳑𘰕 𘲽𘰥𘱸𘲦 𘰺𘱄𘭽𘱸 𘮽𘰣 𘭔𘱄𘱯𘱦 𘲚𘰥𘲐 𘰺𘭸𘰷 𘮦𖿤𘲲 𘰕𘲆𘮒 𘲵𘯺𘬾𘯺 𘬥𘰣𘱮 𘯤 𘲜𘰣𘱤 𘯥 𘯥𘲆𘬷 𘲲𘰓𘰽 𘮧𘬜𘱄 𘰣𘭺 𘮧𘲚 𘲽𘱦 𘯥 𘭃𘭺 𘮽𘲓𘱦 𘰣 𘬘𘭞}


% ==========================================================================
\section{East Asian Writing and CJK Family Selection}
% ==========================================================================

\Term[sort=CJK family selection]{CJK family selection}
[continuous text, shared Han ownership, and explicit regional overrides]
distinguishes the document CJK family from language-specific local families.
Unicode codepoints alone cannot reliably determine the intended regional form
of shared Han characters, while several East Asian writing systems also
require shaping behavior that extends beyond ordinary Han glyph selection.

The specimens in this section therefore use normal-width continuous text
rather than isolated character strings. They exercise paragraph composition,
line breaking, punctuation, regional glyph forms, Hangul composition,
historical jamo, supplementary-plane ideographs, and mixed-script runs.
A compact table at the end isolates shared Han codepoints for direct visual
comparison.


% --------------------------------------------------------------------------
\subsection{Chinese}
% --------------------------------------------------------------------------

\paragraph{Mainland Simplified Chinese --- Noto SC}

This paragraph uses the explicit Simplified Chinese family throughout.

\SC{现代中文排版并不只是把一个个汉字依次放进页面。一个真正的正文环境还需要同时处理标点、数字、书名号、引号、括号以及连续句子之间的换行关系，并且在较长的段落中保持稳定的字面和行距。对于共享的汉字码位，不同地区使用的字体还可能给出不同的字形，因此同一个“骨、直、令、具、角”在不同字体家族中会表现出细微而稳定的区域差异。这里使用一整段连续的简体中文来检查大陆简体字形、标点悬挂附近的布局、正常段落换行以及本地字体选择在长文本中是否能够始终保持有效，而不是只验证几个孤立字符能不能显示。}


\paragraph{Mainland Traditional Chinese --- Noto SC}

Traditional characters can also be deliberately typeset through the Mainland
Chinese family. This separates character repertoire from regional glyph
selection: traditional codepoints do not by themselves imply a Taiwan or
Hong Kong font.

\SC{現代中文排版並不只是把一個個漢字依次放進頁面。一個真正的正文環境還需要同時處理標點、數字、書名號、引號、括號以及連續句子之間的換行關係，並且在較長的段落中保持穩定的字面和行距。對於共享的漢字碼位，不同地區使用的字體還可能給出不同的字形，因此同一個「骨、直、令、具、角」在不同字體家族中會表現出細微而穩定的區域差異。這裡特意使用傳統漢字而仍然選擇大陸字體家族，以檢查字符繁簡與地域字形選擇能否保持彼此獨立，並觀察長段落中共享漢字、繁體字和全角標點的連續排版。}


\paragraph{Taiwan Traditional Chinese --- Noto TC}

The Taiwan family is tested with continuous Taiwan Traditional Chinese rather
than with a short regional-glyph sample.

\TC{正體中文的正文排版除了字符本身之外，也牽涉字型、標點、行距與段落節奏。當文件同時包含不同語言或不同地區的漢字時，程式不能僅依 Unicode 碼位猜測應採用哪一套地區字形，而應由明確的語言或字型選擇來決定。這個段落使用臺灣正體字型，並刻意包含「骨、直、令、具、角、門、國、學、龍」等適合觀察地區字形差異的共享漢字，同時加入引號、頓號、逗號與句號，藉此確認較長的正體中文內容在正常欄寬中能自然換行，而且從段首到段末始終維持同一個本地 CJK family。}


\paragraph{Document Traditional Chinese --- Shanggu}

The following paragraph is deliberately left completely unwrapped. It
therefore remains under the document CJK family, Shanggu, and provides the
control against which the explicit Chinese, Japanese, and Korean families can
be compared.

漢字進入現代數位排版之後，同一個編碼往往可以由許多不同的字體呈現，而字形的歷史層次、地域習慣和設計風格並不會因此消失。傳統中文尤其適合用來觀察這種情形：有些字在不同地區的正文中結構略異，有些字的筆形和部件位置不同，而另一些差別只有在一整段文字的灰度與節奏中才會顯現。這個段落不使用任何語言專用命令，使所有漢字直接落入文件的 Shanggu CJK family；因此它不只是檢查單字是否存在，也檢查文件級 CJK 選擇在較長的連續文本、標點和自動換行之間能否穩定維持。


% --------------------------------------------------------------------------
\subsection{Korean and Historical Hangul}
% --------------------------------------------------------------------------

\paragraph{Modern Korean --- Hangul}

The first Korean specimen deliberately avoids Hanja. It tests ordinary modern
Hangul composition, punctuation, word spacing, and paragraph wrapping.

\KR{한글로 쓰인 긴 문장은 여러 음절 블록이 끊임없이 이어지면서도 낱말 사이의 간격과 문장 부호의 위치가 자연스럽게 유지되어야 한다. 짧은 보기에서는 글꼴이 몇 글자를 제대로 보여 주는지만 확인할 수 있지만, 실제 문단에서는 줄 끝에서 낱말이 어떻게 배치되는지, 괄호와 따옴표가 앞뒤 글자와 어떻게 어울리는지, 여러 줄을 지나도 같은 글꼴 선택이 계속 유지되는지를 함께 살펴볼 수 있다. 이 문단은 완성형 한글이 이어지는 일반적인 한국어 본문을 충분히 길게 배치하여 문장 전체의 모양과 줄 간격을 확인하고, 뒤에 나오는 한자 혼용문과 옛한글 자료를 비교하기 위한 현대 한글 기준 표본으로 사용한다.}


\paragraph{Modern Korean --- mixed Hanja and Hangul}

The next paragraph keeps Hangul and Hanja inside a single explicit Korean
family. Shared Han characters must therefore retain Korean regional forms
without interrupting the surrounding Hangul text.

\KR{韓國語의 現代 文體는 오늘날 대체로 한글만으로 적지만, 歷史ㆍ言語ㆍ哲學ㆍ法律과 같은 분야에서는 漢字를 함께 적은 문헌을 어렵지 않게 만날 수 있다. 이런 漢字混用文을 조판할 때에는 한글 音節과 句讀點 사이에 들어가는 漢字가 갑자기 다른 地域 字形으로 바뀌지 않아야 하며, 韓國ㆍ朝鮮ㆍ漢字ㆍ文化ㆍ文字와 같은 어휘도 앞뒤의 한글 문맥과 한 family 안에서 자연스럽게 이어져야 한다. 이 문단은 國語의 읽기 흐름을 유지하면서 여러 漢字語와 고유어를 섞어, 실제 문헌에서 볼 수 있는 긴 혼용문이 줄바꿈을 거쳐도 안정적으로 조판되는지를 시험한다.}


\paragraph{Jeju language --- corpus specimen}

Jeju is useful here because practical orthographies make active use of the
arae-a and decomposed Hangul sequences that are rare in ordinary modern
Korean. The following is a composite regression specimen assembled from
short attested Jeju-language passages rather than a newly invented
``standard'' Jeju paragraph.

\KR{제주말(濟州말)은 제주섬서 ᄀᆞᆯ아지는 말이다. 제주말모음은 시간 지나멍 여러 가지 변화가 이섯다. 제주섬은 한라산광 ᄌᆞᄁᆞᆺ듸 오름덜록 일뤄진 섬이라. 엿적 화산활동으록 굼부리가 꿰쉉 곤데곤데 ᄋᆢ라 오름이 하영 삿다. 낭(木)은 나무질로 뒌 줄기를 가지고 잇인 여러헤살이 식물이다. 제주말 위키벡과 또시는 제주어 위키백과는 제주말로 운영헴신 위키벡과 언어판이라.}

\paragraph{Middle Korean and Old Hangul --- historical corpus specimen}

The historical specimen deliberately retains decomposed jamo and tone marks.
Rather than normalizing several sources into modern spelling, short genuine
historical fragments are placed consecutively to form a paragraph-level
shaping regression.

\KR{나랏〮말〯ᄊᆞ미〮中듀ᇰ國귁〮에〮 달아〮문ᄍᆞᆼ와〮로서르ᄉᆞᄆᆞᆺ디〮아니〮ᄒᆞᆯᄊᆡ〮　아〮기 ᄇᆡ여〮셔 과ᄀᆞᆯ이〮 가ᄉᆞᆷ〮 알파〮주글〮 ᄃᆞᆺ〮ᄒᆞ〮거든〮　두들겐 開闢ᄒᆞᆯ 젯 므리 疏通ᄒᆞ고 남ᄀᆞᆫ 古今엣 남기 섯것도다　머리〮 알ᄑᆞ고〮 고〮히 막고〮 머리〮와눈〮괘〮 환티〮 아니〮커든〮　긔〮운이〮 우흐〮로 올아〮 머리〮 알포〮미 ᄣᆞ리ᄂᆞᆫ ᄃᆞᆺ〮거든〮　넙고 큰 天造 ᄉᆞ이예 다ᄉᆞᆯ며 어즈러우민ᄃᆞᆯ 엇디 덛덛ᄒᆞᆫ 數ㅣ리오}

The corpus paragraph above tests realistic historical material. The following
synthetic boundary sequence has a different purpose: it explicitly crosses
Hangul Jamo Extended-A, ordinary Jamo, and Extended-B while remaining inside
one shaping run.

\begin{center}
\Huge
\KR{ꥠᅡퟋ〮}
\end{center}

Modern decomposed control:

\begin{center}
\Huge
\KR{각}
\end{center}


% --------------------------------------------------------------------------
\subsection{Japanese and Ryukyuan}
% --------------------------------------------------------------------------

\paragraph{Modern Japanese}

Modern Japanese combines kana, regional Han forms, punctuation, and frequent
script changes inside the same sentence. The entire paragraph remains under
the explicit Japanese family.

\JP{現代日本語の文章では、漢字、ひらがな、カタカナが一つの段落の中で何度も切り替わり、句読点や括弧、かぎ括弧も連続して現れる。短い文字列だけではフォントが表示できることしか分からないが、長い本文にすると、漢字と仮名の大きさの釣り合い、約物の位置、行末での折り返し、そして次の行へ移った後の字形の連続性まで確認できる。この段落では「日本語」「漢字」「文字」「学校」「文化」「骨」「直」「令」「具」「角」などを通常の文脈に入れ、日本語用の地域字形を保ったまま複数行の文章が安定して組版されるかを試している。}


\paragraph{Historical kana orthography}

The opening sentence below is historical text; the continuation is a
showcase-written regression passage that deliberately retains historical kana
spellings such as \JP{ゐる}, \JP{やう}, and \JP{言ふ}. This keeps the
specimen long enough for paragraph testing without presenting newly written
material as an archival transcription.

\JP{真夏の宿屋は空虚であつた。窓の外には白い雲がゆつくり流れ、庭の木立では蟬が絶えず鳴いてゐる。旅人は古い椅子に腰をおろし、まだ來ない馬車の音を待ちながら、時計の針ばかりを眺めてゐた。誰も何も言はないので、午後の光だけが廊下を少しづつ移つて行くやうに見えた。かうして長い時間が過ぎても、玄關の戸は一度も開かず、遠くの道から聞えて來る物音だけが、時々靜かな宿の空氣を動かしてゐるのであつた。}

\paragraph{Japanese Kanbun}

The first sentence is taken from Rai San'yō's \emph{Nihon Gaishi}; the
remainder is a showcase-written Kanbun-style regression continuation. The
purpose here is typographic: Japanese Han glyphs must remain selected even
when no kana occur in the run.

\JP{外史氏曰、吾讀舊志、見鳥羽帝時數下制符。蓋日本文章之傳、漢籍之學與國史之記相因而久矣。昔之學者讀經史、作序記、論人物、述治亂、多用漢文而各成其體。今取其字列於一段、所以驗日本字形之連續、句讀之位置、長行之折返、以及共享漢字在無假名介入時亦不離其本地字體也。文字雖同出一碼、而骨直令具角門國學龍諸字之形、隨地域字體而微異，故尤宜以連續文章觀之。}


\paragraph{Okinawan --- corpus specimen}

Okinawan is represented by a composite specimen from the Wikimedia Okinawan
test project. Its mixed use of kana and Han characters makes it a useful
additional Japanese-family regression, while the orthography should not be
read as a claim that one universal Okinawan spelling standard exists.

\JP{沖縄口ぇー琉球諸語ぬてぃ壱ちやいびーん。うちなーぐちぬ首里那覇方言ぉー琉球国ぬじょーま言葉やびたん。琉球王國や、昔ジジ中心んかい存在さる王国やたん。王や首里天加那志とぅ言ん。首里城や、沖縄県ぬ城やん。琉球王朝ぬ王城っし、沖縄県ぬいっちん大さるぐしくやん。漢字ぇー、大昔唐ぬ国んかいそーじー現りーる文字ぬてぃーちやん。やまとぅ、高麗、るーちゅー、ベトナムぉー、大昔唐ぬ国から漢字伝わてぃちゃん。}

\paragraph{Ainu}

Ainu written in the modern katakana orthography makes extensive use of
Katakana Phonetic Extensions for syllable-final consonants. The following
specimen therefore exercises ordinary katakana together with characters such
as \JP{ㇰ}, \JP{ㇱ}, \JP{ㇺ}, \JP{ㇻ}, and the combining form \JP{ㇷ゚} inside
continuous text.

\JP{イタㇰプリカ　ソモネコㇿカ　シサムモシㇼ　モシㇼソカワ　チヌㇺケニㇱパ　チヌㇺケ　カッケマッ　ウタペラリワ　オカウㇱケタ　クニネネワ　アイヌイタㇰアニ　クイタㇰルウェ　ネワネヤクン　ラモッシワノ　クヤイライケㇷ゚　ネルウェタパンナ。エエパキタ　カニアナㇰネ　アイヌモシㇼ　シシㇼムカ　ニㇷ゚タニコタン　コアパマカ　クネルウェネ　ウッチケクニㇷ゚　ラカサㇰペ。}

% --------------------------------------------------------------------------
\subsection{Vietnamese}
% --------------------------------------------------------------------------

The same constitutional passage is used throughout this subsection so that
the specimens differ primarily in writing system or font selection rather
than in linguistic content. Quốc Ngữ is first typeset with the document Latin
family, Libertinus, and then with the dedicated Vietnamese family. The
corresponding Hán-Nôm text uses the VinaWiki repertoire.

\paragraph{Quốc Ngữ --- document Latin family}

The first specimen is deliberately unwrapped and therefore remains in the
document Latin family, Libertinus. It tests Vietnamese diacritics without a
language-specific font override.

Trải qua mấy nghìn năm lịch sử, Nhân dân Việt Nam lao động cần cù, sáng tạo, đấu tranh anh dũng để dựng nước và giữ nước, đã hun đúc nên truyền thống yêu nước, đoàn kết, nhân nghĩa, kiên cường, bất khuất và xây dựng nên nền văn hiến Việt Nam. Từ năm 1930, dưới sự lãnh đạo của Đảng Cộng sản Việt Nam do Chủ tịch Hồ Chí Minh sáng lập và rèn luyện, Nhân dân ta tiến hành cuộc đấu tranh lâu dài, đầy gian khổ, hy sinh vì độc lập, tự do của dân tộc, vì hạnh phúc của Nhân dân. Cách mạng tháng Tám thành công, ngày 2 tháng 9 năm 1945, Chủ tịch Hồ Chí Minh đọc Tuyên ngôn độc lập, khai sinh ra nước Việt Nam dân chủ cộng hòa, nay là Cộng hòa xã hội chủ nghĩa Việt Nam.


\paragraph{Quốc Ngữ --- Vietnamese family}

The same text is repeated through the explicit Vietnamese local family,
allowing the dedicated Vietnamese font to be compared directly with
Libertinus under identical paragraph conditions.

\VI{Trải qua mấy nghìn năm lịch sử, Nhân dân Việt Nam lao động cần cù, sáng tạo, đấu tranh anh dũng để dựng nước và giữ nước, đã hun đúc nên truyền thống yêu nước, đoàn kết, nhân nghĩa, kiên cường, bất khuất và xây dựng nên nền văn hiến Việt Nam. Từ năm 1930, dưới sự lãnh đạo của Đảng Cộng sản Việt Nam do Chủ tịch Hồ Chí Minh sáng lập và rèn luyện, Nhân dân ta tiến hành cuộc đấu tranh lâu dài, đầy gian khổ, hy sinh vì độc lập, tự do của dân tộc, vì hạnh phúc của Nhân dân. Cách mạng tháng Tám thành công, ngày 2 tháng 9 năm 1945, Chủ tịch Hồ Chí Minh đọc Tuyên ngôn độc lập, khai sinh ra nước Việt Nam dân chủ cộng hòa, nay là Cộng hòa xã hội chủ nghĩa Việt Nam.}


\paragraph{Hán-Nôm --- VinaWiki repertoire}

The Hán-Nôm specimen corresponds to the Quốc Ngữ passage above rather than
serving as an unrelated coverage string. Because the local family uses the
VinaWiki font, the specimen is taken from the VinaWiki transcription itself
and therefore tests the repertoire for which that font is intended.

\HN{𱱇過𠇍𠦳𢆥歷史、人民越南勞動勤劬、創造、鬥爭英勇抵𥩯渃吧𡨺渃、㐌熏𤒘𢧚傳統𢞅渃、團結、仁義、堅強、不屈吧𡏦𥩯𢧚𡋂文獻越南。自𢆥1930、𨑜事領導𧵑黨共產越南由主席胡志明創立吧㷙鍊、人民些進行局鬥爭𥹰󠄁𨱽、𣹓艱苦、犧牲爲獨立、自由𧵑民族、爲幸福𧵑人民。革命𣎃𠔭成功、𣈜2𣎃9𢆥1945、主席胡志明讀宣言獨立、開生𫥨渃越南民主共和、𫢩羅共和社會主義越南。}


% --------------------------------------------------------------------------
\subsection{Tangut}
% --------------------------------------------------------------------------

Wikisource provides a sizeable Unicode Tangut corpus, including poetry,
Buddhist literature, maxims, and the Tangut text of the Liángzhōu bilingual
stele. The following paragraph is therefore a composite corpus specimen:
several short attested sequences from independent Tangut texts are joined
only for visual regression. It is not presented as one continuous historical
work.

\TG{𘘞𗳝𘛮𗫉𗥦𗤄，𗼑𗼑𗫉𗰜𗰓𗨛，𘚜𗨤𗤮𘁾𗣃𗅋𘃪，𗼇𗗙𘋥𘝿𘎪，𗵒𗹶𗲈𗃵𘓺𘋨𗶠，𗦵𘇇𗱢𗆐𗲌𗽓𘏋，𗌜𗎫𗠒𗂥𗅋𗚛𘓁𘟂𗫶，𘚶𗢵𘄩𗘴𘜑𗃱𗀸𗀸𗏹𗅋𗁿，𗒘𘓷𗰜𘕿𗅋𗯗𘓁𘟂𗫶，𗌮𗈦𗪙𘝇𗤁𗵘𘆗𘚢，𘈒𗤄𗢳𘒣𗎫𘟀𗢳𗵆，𘘚𘒣𗰜𗎫𗮔𗩱𗫂𘟀𘟂，𗮔𗰜𗋕𗡶𘉐𗭍𗅋𗋃𗫂𗎫𘟂。}


% --------------------------------------------------------------------------
\subsection{WenJin Multi-Plane Coverage}
% --------------------------------------------------------------------------

WenJin is different from the language-specific families above. It is exposed
as one public local family whose internal fallback configuration spans
multiple CJK planes. The paragraph below is therefore intentionally a
coverage stress test rather than a linguistic specimen: ordinary Han text is
interleaved with supplementary-plane characters while the public command
remains unchanged.

\WJ{中國文字的編碼範圍遠遠超過基本多文種平面，因此只測試一二三天地玄黃這類常用字並不足以確認大型漢字字體的覆蓋是否正常。這一段先由普通漢字開始，隨後插入第二平面的𠀀、𠮷、𡈽、𡌛、𣎴，再回到典籍文字與古今學術，之後繼續進入第三平面的𰀀、𰀁、𰀂、𰀃、𰀄。當不同平面的字符在同一個長段落中交錯出現時，使用者仍然只看見一個 \texttt{\textbackslash WJ} family；Plane 0、Plane 2 和 Plane 3 的具體 fallback 應由內部配置完成，而不應要求正文作者為不同 Unicode plane 手動切換命令。}

The following table complements the paragraph test with deterministic block
coverage samples. For each CJK Unified Ideographs block, the first twenty
encoded ideographs are shown. This makes missing coverage or an incorrect
Plane 0, Plane 2, or Plane 3 fallback immediately visible and keeps the
regression sample stable across releases.

\begin{TableBookX}
  {L L L}
  {WenJin CJK block coverage}
  {tab:wenjin-blocks}

Block & Range & Twenty-character specimen \\
\midrule

CJK Unified Ideographs
& U+4E00--U+9FFF
& \WJ{一丁丂七丄丅丆万丈三上下丌不与丏丐丑丒专} \\

Extension A
& U+3400--U+4DBF
& \WJ{㐀㐁㐂㐃㐄㐅㐆㐇㐈㐉㐊㐋㐌㐍㐎㐏㐐㐑㐒㐓} \\

Extension B
& U+20000--U+2A6DF
& \WJ{𠀀𠀁𠀂𠀃𠀄𠀅𠀆𠀇𠀈𠀉𠀊𠀋𠀌𠀍𠀎𠀏𠀐𠀑𠀒𠀓} \\

Extension C
& U+2A700--U+2B73F
& \WJ{𪜀𪜁𪜂𪜃𪜄𪜅𪜆𪜇𪜈𪜉𪜊𪜋𪜌𪜍𪜎𪜏𪜐𪜑𪜒𪜓} \\

Extension D
& U+2B740--U+2B81F
& \WJ{𫝀𫝁𫝂𫝃𫝄𫝅𫝆𫝇𫝈𫝉𫝊𫝋𫝌𫝍𫝎𫝏𫝐𫝑𫝒𫝓} \\

Extension E
& U+2B820--U+2CEAF
& \WJ{𫠠𫠡𫠢𫠣𫠤𫠥𫠦𫠧𫠨𫠩𫠪𫠫𫠬𫠭𫠮𫠯𫠰𫠱𫠲𫠳} \\

Extension F
& U+2CEB0--U+2EBEF
& \WJ{𬺰𬺱𬺲𬺳𬺴𬺵𬺶𬺷𬺸𬺹𬺺𬺻𬺼𬺽𬺾𬺿𬻀𬻁𬻂𬻃} \\

Extension I
& U+2EBF0--U+2EE5F
& \WJ{𮯰𮯱𮯲𮯳𮯴𮯵𮯶𮯷𮯸𮯹𮯺𮯻𮯼𮯽𮯾𮯿𮰀𮰁𮰂𮰃} \\

Extension G
& U+30000--U+3134F
& \WJ{𰀀𰀁𰀂𰀃𰀄𰀅𰀆𰀇𰀈𰀉𰀊𰀋𰀌𰀍𰀎𰀏𰀐𰀑𰀒𰀓} \\

Extension H
& U+31350--U+323AF
& \WJ{𱍐𱍑𱍒𱍓𱍔𱍕𱍖𱍗𱍘𱍙𱍚𱍛𱍜𱍝𱍞𱍟𱍠𱍡𱍢𱍣} \\

Extension J
& U+323B0--U+3347F
& \WJ{𲎰𲎱𲎲𲎳𲎴𲎵𲎶𲎷𲎸𲎹𲎺𲎻𲎼𲎽𲎾𲎿𲏀𲏁𲏂𲏃} \\

\end{TableBookX}


% --------------------------------------------------------------------------
\subsection{Regional Forms of Shared Han Characters}
% --------------------------------------------------------------------------

The preceding specimens test normal paragraphs. Only here are shared
codepoints deliberately isolated. Every cell below contains the same Unicode
character; the visible difference should therefore come from the selected
regional family rather than from different encoded variants.

\begin{TableBookX}
  {L C C C C C}
  {Regional forms of shared Han codepoints}
  {tab:cjk-regional-forms}

Character & Shanggu & Noto SC & Noto TC & Japanese & Korean \\
\midrule

骨 & 骨 & \SC{骨} & \TC{骨} & \JP{骨} & \KR{骨} \\

直 & 直 & \SC{直} & \TC{直} & \JP{直} & \KR{直} \\

令 & 令 & \SC{令} & \TC{令} & \JP{令} & \KR{令} \\

具 & 具 & \SC{具} & \TC{具} & \JP{具} & \KR{具} \\

角 & 角 & \SC{角} & \TC{角} & \JP{角} & \KR{角} \\

門 & 門 & \SC{門} & \TC{門} & \JP{門} & \KR{門} \\

國 & 國 & \SC{國} & \TC{國} & \JP{國} & \KR{國} \\

學 & 學 & \SC{學} & \TC{學} & \JP{學} & \KR{學} \\

龍 & 龍 & \SC{龍} & \TC{龍} & \JP{龍} & \KR{龍} \\

\end{TableBookX}

% ==========================================================================
\section{Document Features}
% ==========================================================================

\Term[sort=feature system]{feature system}
[composable document modules] is demonstrated throughout this PDF, but this
section gives each major public feature an explicit specimen as well. The
showcase therefore tests not only isolated modules but also whether they
coexist in a realistic document.

% --------------------------------------------------------------------------
\subsection{Mathematics}
% --------------------------------------------------------------------------

Mathematical font selection is deliberately left at IMPE's default setting.
The following single display combines several common mathematical alphabets,
operators, delimiters, accents, matrices, integrals, sums, limits, radicals,
and nested scripts as a compact stress specimen.

\begingroup
\small
\[
\begin{aligned}
\mathcal Z_{\beta,\mu}
&\coloneqq
\operatorname{Tr}_{\mathscr H}
\exp\!\left[
-\beta
\left(
\widehat{\mathbf H}
-\mu\widehat{\mathbf N}
\right)
\right]
\\[0.5em]
&=
\sum_{n=0}^{\infty}
\frac{1}{n!}
\int_{\mathbb R^{3n}}
\prod_{j=1}^{n}
\frac{d^3\mathbf p_j}{(2\pi\hbar)^3}
\,
\exp\!\left[
-\sum_{j=1}^{n}
\sqrt{
m_j^2c^4+c^2\lVert\mathbf p_j\rVert^2
}
\right]
\\[0.7em]
&\quad\times
\det
\begin{pmatrix}
\displaystyle
\frac{\partial^2\mathcal L}
     {\partial q^i\partial q^j}
&
\displaystyle
\nabla_\mu F^{\mu\nu}
&
\displaystyle
\left\langle
\psi
\middle|
\widehat{\mathbf H}
\middle|
\phi
\right\rangle
\\[0.9em]
\displaystyle
\sum_{k=1}^{\infty}\frac{(-1)^{k+1}}{k^2}
&
\displaystyle
\sqrt[3]{1+\sqrt{1+\lambda^2}}
&
\displaystyle
\frac{\mathbb E[X^4]-3\mathbb E[X^2]^2}
     {\operatorname{Var}(X)^2}
\\[0.9em]
\displaystyle
\oint_{|z|=1}
\frac{\log(1+z)}{z^{r+1}}\,dz
&
\displaystyle
\prod_{p\in\mathbb P}
\left(1-p^{-s}\right)^{-1}
&
\displaystyle
\lim_{\varepsilon\to0^+}
\frac{f(x+\varepsilon)-f(x-\varepsilon)}
     {2\varepsilon}
\end{pmatrix}
\\[1em]
&\quad+
\underbrace{
\widetilde{\mathbf A}^{\,\dagger}{}_{\mu}{}^{\nu}
\,
\overline{\mathbf B}_{\nu}{}^{\rho}
\,
\widehat{\mathbf C}_{\rho}{}^{\mu}
}_{\text{accented tensor contraction}}
+
\frac{
\Gamma\!\left(\frac12+i\alpha\right)
\zeta(2s)
}{
\sqrt{2\pi\sigma^2}
}
\\[0.8em]
&\quad+
\begin{cases}
\displaystyle
e^{-1/x^2}\sin(\pi/x),
& x\ne0,
\\[0.8em]
0,
& x=0,
\end{cases}
\\[0.8em]
&\quad
\xrightarrow[\alpha\to0]{\beta\to\infty}
\left(
\bigoplus_{j=1}^{r}\mathbb C^{\,n_j}
\right)
\otimes
\bigwedge\nolimits^{k}V
\\[0.8em]
&\quad
\Longrightarrow
\mathbb P
\left(
|X-\mathbb E[X]|\ge t
\right)
\le
\frac{\operatorname{Var}(X)}{t^2},
\qquad
t>0,
\\[0.8em]
&\quad
\alpha,\beta,\gamma,\delta,\varepsilon,\theta,\vartheta,\lambda,\mu,\nu,
\xi,\pi,\varpi,\rho,\varrho,\sigma,\phi,\varphi,\chi,\psi,\omega
\in\mathbb R .
\end{aligned}
\]
\endgroup

\subsection{Tables}

Table~\ref{tab:feature-summary} uses the IMPE book-table environment rather
than a raw \texttt{tabular} wrapper.

\begin{TableBookX}
  {L C C}
  {Feature integration summary}
  {tab:feature-summary}
Feature & Public route & Exercised here \\
\midrule
Mathematics & \texttt{math} & yes \\
Tables & \texttt{tables} & yes \\
Images & \texttt{image} & yes \\
Citations & \texttt{citations} & yes \\
Index & \texttt{index} & yes \\
Hyperlinks & \texttt{hyperlinks} & yes \\
Running headers & \texttt{headers} & yes \\
List environments & \texttt{lists\_envs} & yes \\
\end{TableBookX}

\subsection{Lists and Example Blocks}

A normal list remains ordinary \LaTeX:

\begin{enumerate}
  \item automatic font ownership;
  \item script-specific behavior where required;
  \item reusable document features;
  \item stable public interfaces.
\end{enumerate}

The IMPE \texttt{ExampleBlock} environment provides a distinct indented
specimen:

\begin{ExampleBlock}
A canonical showcase should be readable as a document, useful as a visual
regression specimen, and representative of the public interface that new
documents are expected to use.
\end{ExampleBlock}

\subsection{Images}

The image feature is exercised with the standard \TeX~distribution's example
image when it is available. This keeps the canonical source self-contained
without committing an otherwise meaningless binary illustration.

\IfFileExists{example-image-a.pdf}{%
  \begin{OneImage}
    [htbp]
    [0.55\linewidth]
    [0.32\textheight]
    {example-image-a}
    [Image-feature specimen using the standard \TeX~example image.]
    [fig:showcase-image]
  \end{OneImage}

  Figure~\ref{fig:showcase-image} also exercises cross-reference hyperlinks.
}{%
  \begin{center}
    \fbox{%
      \parbox{0.62\linewidth}{%
        \centering
        \vspace{2em}
        \textbf{Image specimen unavailable}\\[0.5em]
        \texttt{example-image-a.pdf} was not found in this \TeX~installation.
        \vspace{2em}
      }%
    }
  \end{center}
}

\subsection{Citations and Bibliography}

Citation support is loaded through the feature system and configured before
feature activation. A normal textual citation can therefore be written as
\textcite{knuth1984texbook}, while a parenthetical citation appears as
\parencite{knuth1984texbook}.

The bibliography is printed near the end of the document.

\subsection{Index Terms}

Several concepts in earlier sections were inserted with the public
\texttt{\textbackslash Term} interface. For example,
\Term[sort=canonical showcase]{canonical showcase}
[release-time specimen] itself is now indexed. The generated index at the end
of the PDF therefore also exercises indexed terms, optional descriptions, and
hyperlinks back to their first occurrence.

\subsection{Hyperlinks, Footnotes, and Running Headers}

The table of contents is clickable, numbered headings carry stable targets,
and section headings can link back toward their table-of-contents entries.
The running header above this page is produced by the \texttt{headers}
feature.

Footnote markers are also bidirectional.\footnote{Clicking the footnote marker
jumps to this note; the note marker can link back to the original position.}

For an external-link specimen, the project repository can be represented with
a normal hyperlink:
\href{https://github.com/KelvinYangBK67/IMPE-LaTeX-System}
{IMPE \LaTeX~System on GitHub}.


\printbibliography[
  heading=bibintoc,
  title={References}
]

\printindex

\end{document}
